\documentclass[10pt]{beamer}

% --- Modern Theme ---
\usetheme{metropolis}
\usecolortheme{default}
\setbeamercolor{background canvas}{bg=white}

% --- Packages ---
\usepackage[utf8]{inputenc}
\usepackage{booktabs}
\usepackage{array}
\usepackage{tabularx}
\usepackage{amsmath}
\usepackage{amssymb}
\usepackage[scale=2]{ccicons}
\usepackage{pgfplots}
\usepgfplotslibrary{dateplot}
\usepackage{xspace}

% --- Title Information ---
\title{Foundations of Artificial Intelligence}
\subtitle{Part V: Natural Language Processing (Introductory)}
\institute{Department of Computer Science}
\date{}

\begin{document}

\maketitle

\begin{frame}{Lecture Roadmap}
    \setbeamertemplate{section in toc}[sections numbered]
    \tableofcontents[hideallsubsections]
\end{frame}

\begin{frame}{Learning Objectives}
    By the end of this lecture, you should be able to:
    \begin{itemize}
        \item Apply basic text preprocessing techniques: tokenization, lowercasing, stopword removal, stemming.
        \item Build and interpret Bag-of-Words representations for text.
        \item Implement basic sentiment analysis using word-level patterns.
        \item Discuss real-world NLP applications: chatbots, translation, and speech systems.
        \item Recognize the role of language models and the limitations of simple approaches.
    \end{itemize}
\end{frame}

\section{Text Preprocessing}

\begin{frame}{Why Preprocessing Matters}
    Raw text is messy:
    \begin{itemize}
        \item Mixed case, punctuation, special characters.
        \item Irrelevant words (``the'', ``a'', ``is'').
        \item Different forms of the same word (``run'', ``runs'', ``running'').
    \end{itemize}
    \vspace{0.2cm}
    Preprocessing standardizes text so that algorithms see meaningful units consistently.

    \vspace{0.2cm}
    \textbf{Core steps:}
    \begin{enumerate}
        \item Tokenization
        \item Lowercasing and normalization
        \item Stopword removal
        \item Stemming/lemmatization
    \end{enumerate}
\end{frame}

\begin{frame}{Tokenization}
    \textbf{Goal:} Split text into individual words (tokens).

    \vspace{0.2cm}
    \textbf{Simple approach:} split on whitespace and punctuation.
    \begin{itemize}
        \item Input: \texttt{``Hello, world! How are you?''}
        \item Output: \texttt{[``Hello'', ``world'', ``How'', ``are'', ``you'']}
    \end{itemize}

    \vspace{0.2cm}
    \textbf{Challenges:}
    \begin{itemize}
        \item Contractions: \texttt{``don't''} $\to$ keep as one or split?
        \item Hyphenated words: \texttt{``state-of-the-art''}.
        \item URLs, emails, numbers: preserve or normalize?
    \end{itemize}

    \vspace{0.2cm}
    Most tools use regex-based or language-aware tokenizers.
\end{frame}

\begin{frame}{Lowercasing and Normalization}
    \textbf{Lowercasing:} Convert all text to lowercase.
    \begin{itemize}
        \item Reduces vocabulary size.
        \item Treats \texttt{``Hello''}, \texttt{``hello''}, \texttt{``HELLO''} identically.
    \end{itemize}

    \vspace{0.2cm}
    \textbf{Other normalizations:}
    \begin{itemize}
        \item Remove accents: \texttt{``café''} $\to$ \texttt{``cafe''}.
        \item Replace numbers with placeholders: \texttt{``123 Main St''} $\to$ \texttt{``[NUM] Main St''}.
        \item Expand contractions: \texttt{``don't''} $\to$ \texttt{``do not''}.
    \end{itemize}

    \vspace{0.2cm}
    Trade-off: lose information (e.g., proper nouns) for uniformity.
\end{frame}

\begin{frame}{Stopword Removal}
    \textbf{Stopwords:} Common words with little discriminative power.
    \begin{itemize}
        \item Articles: ``a'', ``an'', ``the''
        \item Pronouns: ``I'', ``you'', ``he'', ``she''
        \item Prepositions: ``in'', ``on'', ``at'', ``to''
        \item Conjunctions: ``and'', ``or'', ``but''
    \end{itemize}

    \vspace{0.2cm}
    \textbf{Example:}
    \begin{itemize}
        \item Before: ``The cat is on the mat and it is very soft.''
        \item After: ``cat mat soft''
    \end{itemize}

    \vspace{0.2cm}
    \textbf{Benefit:} Reduces noise and speeds up processing.

    \vspace{0.1cm}
    \textbf{Caveat:} Domain-dependent; task may require keeping some stopwords.
\end{frame}

\begin{frame}{Stemming and Lemmatization}
    \textbf{Stemming:} Remove suffixes to get word root.
    \begin{itemize}
        \item Input: ``running'', ``runs'', ``ran'', ``runner''
        \item Output (Porter Stemmer): ``run'', ``run'', ``ran'', ``runner''
        \item Heuristic-based; may produce non-words.
    \end{itemize}

    \vspace{0.2cm}
    \textbf{Lemmatization:} Reduce to canonical dictionary form (lemma).
    \begin{itemize}
        \item Input: ``running'', ``runs'', ``ran'', ``runner''
        \item Output: ``run'', ``run'', ``go'', ``run''
        \item Requires linguistic knowledge; more accurate but slower.
    \end{itemize}

    \vspace{0.2cm}
    \textbf{Trade-off:} Stemming is fast but crude; lemmatization is precise but needs external resources.
\end{frame}

\begin{frame}{Preprocessing Pipeline Example}
    \small
    \setlength{\tabcolsep}{8pt}
    \renewcommand{\arraystretch}{1.2}
    \begin{tabular}{>{\raggedright\arraybackslash}p{3cm} >{\raggedright\arraybackslash}p{5.5cm}}
        \toprule
        \textbf{Step} & \textbf{Result} \\
        \midrule
        Original & ``The analysis of running cats and dogs is fascinating!'' \\
        Tokenize & [``The'', ``analysis'', ``of'', ``running'', ``cats'', ``and'', ``dogs'', ``is'', ``fascinating'', ``!''] \\
        Lowercase & [``the'', ``analysis'', ``of'', ``running'', ``cats'', ``and'', ``dogs'', ``is'', ``fascinating'', ``!''] \\
        Remove punct. & [``the'', ``analysis'', ``of'', ``running'', ``cats'', ``and'', ``dogs'', ``is'', ``fascinating''] \\
        Remove stopwords & [``analysis'', ``running'', ``cats'', ``dogs'', ``fascinating''] \\
        Stem & [``analysi'', ``run'', ``cat'', ``dog'', ``fascin''] \\
        \bottomrule
    \end{tabular}
\end{frame}

\section{Bag-of-Words Representation}

\begin{frame}{From Tokens to Vectors}
    \textbf{Bag-of-Words (BoW):} Represent text as unordered collection of word counts.

    \vspace{0.2cm}
    Key idea: ignore order, focus on word frequency.

    \vspace{0.3cm}
    \textbf{Example:}
    \begin{itemize}
        \item Document 1: ``I love machine learning''
        \item Document 2: ``I love cats and dogs''
        \item Vocabulary: \{I, love, machine, learning, cats, dogs, and\}
    \end{itemize}

    \vspace{0.2cm}
    BoW vectors (counts):
    \begin{itemize}
        \item Doc 1: $[1, 1, 1, 1, 0, 0, 0]$
        \item Doc 2: $[1, 1, 0, 0, 1, 1, 1]$
    \end{itemize}
\end{frame}

\begin{frame}{BoW Variants: TF and TF-IDF}
    \textbf{Term Frequency (TF):} Count of word in document.
    \[
        \text{TF}(w, d) = \text{count}(w, d)
    \]

    \vspace{0.2cm}
    \textbf{Problem:} Common words dominate. Longer documents have higher counts.

    \vspace{0.25cm}
    \textbf{TF-IDF:} Weight by inverse document frequency.
    \[
        \text{TF-IDF}(w, d) = \text{TF}(w, d) \times \log\left(\frac{N}{|\{d : w \in d\}|}\right)
    \]
    where $N$ is total documents and denominator counts how many documents contain $w$.

    \vspace{0.2cm}
    \textbf{Intuition:} Rare words in a document get high weight.
\end{frame}

\begin{frame}{BoW Strengths and Limitations}
    \textbf{Strengths:}
    \begin{itemize}
        \item Simple to implement and understand.
        \item Works well with many classifiers (Naive Bayes, SVM, etc.).
        \item Computationally efficient.
    \end{itemize}

    \vspace{0.2cm}
    \textbf{Limitations:}
    \begin{itemize}
        \item \textbf{Order-independent:} ``dog bites man'' and ``man bites dog'' are identical.
        \item \textbf{Loses context:} Single words may be ambiguous (``bank'': financial or riverside?).
        \item \textbf{Sparse representations:} vocabulary can be very large (high-dimensional).
        \item \textbf{No semantic similarity:} ``dog'' and ``cat'' are unrelated vectors.
    \end{itemize}

    \vspace{0.2cm}
    \textit{Word embeddings (e.g., Word2Vec, GloVe) address some limitations.}
\end{frame}

\begin{frame}{Document Similarity with BoW}
    \textbf{Goal:} Measure how similar two documents are.

    \vspace{0.2cm}
    \textbf{Cosine Similarity:}
    \[
        \cos(\theta) = \frac{\vec{d}_1 \cdot \vec{d}_2}{|\vec{d}_1| \cdot |\vec{d}_2|}
    \]
    where $\vec{d}_1, \vec{d}_2$ are BoW or TF-IDF vectors.

    \vspace{0.2cm}
    \textbf{Range:} $[{-1, 1}]$ (or $[0, 1]$ for non-negative vectors).
    \begin{itemize}
        \item 1 = identical direction (very similar).
        \item 0 = orthogonal (no similarity).
        \item $-1$ = opposite (rare in BoW).
    \end{itemize}

    \vspace{0.2cm}
    \textbf{Application:} Information retrieval, plagiarism detection, document clustering.
\end{frame}

\section{Basic Sentiment Analysis}

\begin{frame}{What Is Sentiment Analysis?}
    \textbf{Goal:} Determine the emotional tone or opinion in text.

    \vspace{0.2cm}
    \textbf{Common task:} Binary classification.
    \begin{itemize}
        \item Positive vs. Negative
        \item Example: product reviews, social media posts.
    \end{itemize}

    \vspace{0.2cm}
    \textbf{More nuanced:} Multi-class (positive, neutral, negative) or regression (score 1--5).

    \vspace{0.2cm}
    \textbf{Why it matters:}
    \begin{itemize}
        \item Brand monitoring.
        \item Customer feedback analysis.
        \item Social listening.
    \end{itemize}
\end{frame}

\begin{frame}{Lexicon-Based Approach}
    \textbf{Idea:} Use a sentiment lexicon (dictionary of words with known polarity).

    \vspace{0.2cm}
    \textbf{Example lexicon:}
    \small
    \begin{tabularx}{0.9\textwidth}{l l}
        \textbf{Positive} & good, great, excellent, love, amazing \\
        \textbf{Negative} & bad, terrible, hate, awful, disappointing
    \end{tabularx}

    \normalsize
    \vspace{0.2cm}
    \textbf{Algorithm:}
    \begin{enumerate}
        \item Tokenize and preprocess text.
        \item Count positive and negative words.
        \item Assign label: positive if more positive, else negative.
    \end{enumerate}

    \vspace{0.2cm}
    \textbf{Example:} ``This movie is great and I loved it!''
    \begin{itemize}
        \item Positive words: great, loved (count = 2).
        \item Negative words: none (count = 0).
        \item Result: Positive.
    \end{itemize}
\end{frame}

\begin{frame}{Challenges in Sentiment Analysis}
    \textbf{Negation:} ``I did \textbf{not} like this movie.''
    \begin{itemize}
        \item ``like'' is positive, but negated.
        \item Simple lexicon approach fails.
    \end{itemize}

    \vspace{0.2cm}
    \textbf{Sarcasm:} ``Oh, great, another traffic jam.''
    \begin{itemize}
        \item ``great'' is positive, but context is negative.
        \item Requires understanding intent.
    \end{itemize}

    \vspace{0.2cm}
    \textbf{Domain-specific language:} ``sick'' can be negative or positive (slang).

    \vspace{0.2cm}
    \textbf{Ambiguous words:} ``interesting'' could be positive or negative depending on context.

    \vspace{0.2cm}
    \textit{Advanced methods: machine learning, neural networks, contextual embeddings.}
\end{frame}

\begin{frame}{Feature-Based Sentiment Analysis}
    Extend BoW with engineered features:
    \begin{itemize}
        \item \textbf{Word polarity:} sentiment score of each word (e.g., AFINN, TextBlob).
        \item \textbf{Negation handling:} flip polarity if preceded by negation.
        \item \textbf{Capitalization:} ``AMAZING'' might indicate strong emotion.
        \item \textbf{Punctuation:} multiple exclamation marks ``!!!'' suggest intensity.
        \item \textbf{Part-of-speech:} adjectives often carry opinion.
    \end{itemize}

    \vspace{0.2cm}
    \textbf{Workflow:}
    \begin{enumerate}
        \item Extract features from preprocessed text.
        \item Train classifier (Naive Bayes, SVM, Logistic Regression).
        \item Predict sentiment on new reviews.
    \end{enumerate}
\end{frame}

\begin{frame}{Sentiment Analysis Performance Metrics}
    \textbf{Classification metrics:}
    \begin{itemize}
        \item \textbf{Accuracy:} fraction of correct predictions.
        \item \textbf{Precision:} of predicted positives, how many are truly positive.
        \item \textbf{Recall:} of truly positive cases, how many are found.
        \item \textbf{F1-score:} harmonic mean of precision and recall.
    \end{itemize}

    \vspace{0.2cm}
    \textbf{Example confusion matrix (binary):}
    \small
    \begin{tabularx}{0.8\textwidth}{l c c}
        \toprule
        & \textbf{Predicted Pos} & \textbf{Predicted Neg} \\
        \midrule
        \textbf{Actual Pos} & TP & FN \\
        \textbf{Actual Neg} & FP & TN \\
        \bottomrule
    \end{tabularx}

    \normalsize
    \vspace{0.2cm}
    Precision $= \frac{\text{TP}}{\text{TP}+\text{FP}}$, \quad Recall $= \frac{\text{TP}}{\text{TP}+\text{FN}}$
\end{frame}

\section{NLP Applications}

\begin{frame}{Real-World NLP Applications}
    NLP powers many systems we use daily:
    \begin{itemize}
        \item \textbf{Chatbots \& Virtual Assistants:} Siri, Alexa, ChatGPT.
        \item \textbf{Machine Translation:} Google Translate, DeepL.
        \item \textbf{Speech Recognition \& Synthesis:} Dictation, text-to-speech.
        \item \textbf{Information Retrieval:} Search engines, document ranking.
        \item \textbf{Named Entity Recognition:} Extract people, places, organizations.
        \item \textbf{Question Answering:} Systems that find answers from text.
        \item \textbf{Text Summarization:} Automatic document or article condensing.
    \end{itemize}

    \vspace{0.2cm}
    \textit{Modern systems often combine NLP with neural networks and large language models.}
\end{frame}

\begin{frame}{Chatbots}
    \textbf{Goal:} Simulate conversation with a human via text (or speech + NLP).

    \vspace{0.2cm}
    \textbf{Architecture levels:}
    \begin{itemize}
        \item \textbf{Rule-based:} Pattern matching, lookup tables. Limited but predictable.
        \item \textbf{ML-based:} Learns from training data. Intent classification + response generation.
        \item \textbf{Neural (Sequence-to-Seq):} Encoder-decoder models, transformer-based (e.g., GPT).
    \end{itemize}

    \vspace{0.2cm}
    \textbf{Challenge:} Understanding context and maintaining conversation flow.

    \vspace{0.2cm}
    \textbf{Current frontier:} Large language models with few-shot learning enable more natural dialogue.
\end{frame}

\begin{frame}{Machine Translation}
    \textbf{Goal:} Translate text from one language to another.

    \vspace{0.2cm}
    \textbf{Historical approaches:}
    \begin{itemize}
        \item \textbf{Rule-based:} Hand-coded grammar rules and bilingual dictionaries.
        \item \textbf{Statistical MT (SMT):} Learn probabilities from aligned parallel texts.
        \item \textbf{Neural MT (NMT):} Encoder-decoder with attention mechanisms.
    \end{itemize}

    \vspace{0.2cm}
    \textbf{Key challenge:} Structural differences between languages, idioms, context.

    \vspace{0.2cm}
    \textbf{Sequence-to-Sequence with Attention:}
    \begin{itemize}
        \item Encoder processes source language.
        \item Decoder generates target language.
        \item Attention focuses on relevant source words.
    \end{itemize}

    \vspace{0.2cm}
    \textit{Modern systems (Google Translate, DeepL) use transformer models.}
\end{frame}

\begin{frame}{Speech Recognition and Synthesis}
    \textbf{Speech Recognition:} Convert audio to text.
    \begin{itemize}
        \item Acoustic model: features from audio $\to$ phonemes.
        \item Language model: sequence of words, penalizes unlikely sentences.
        \item Decoding: find best word sequence given audio.
    \end{itemize}

    \vspace{0.2cm}
    \textbf{Speech Synthesis (TTS):} Convert text to audio.
    \begin{itemize}
        \item Text analysis: parse input, normalize abbreviations.
        \item Linguistic feature extraction: phoneme, duration, pitch.
        \item Vocoder: generate waveform from features (can sound robotic or natural).
    \end{itemize}

    \vspace{0.2cm}
    \textbf{Modern methods:} Neural networks and end-to-end learning improve naturalness and robustness.
\end{frame}

\begin{frame}{NLP Challenges and Limitations}
    \textbf{Fundamental issues:}
    \begin{itemize}
        \item \textbf{Ambiguity:} Words and sentences have multiple meanings.
        \item \textbf{Context dependence:} Meaning relies on surrounding text and world knowledge.
        \item \textbf{Idioms \& Figurative language:} ``raining cats and dogs'' is not literal.
        \item \textbf{Multilingual:} Each language has unique structure and resources.
    \end{itemize}

    \vspace{0.2cm}
    \textbf{Data challenges:}
    \begin{itemize}
        \item Labeled data is expensive to create.
        \item Imbalanced classes (e.g., rare entities).
        \item Domain shift: model trained on news fails on social media.
    \end{itemize}

    \vspace{0.2cm}
    \textit{Large Language Models (LLMs) have dramatically improved performance through scale and pretraining.}
\end{frame}

\begin{frame}{The Role of Language Models}
    \textbf{Language models:} Predict next word given context.
    \[
        P(\text{word}_t | \text{word}_{1:t-1})
    \]

    \vspace{0.2cm}
    \textbf{Impact:}
    \begin{itemize}
        \item Foundation for many NLP tasks (translation, question answering).
        \item Transfer learning: pretrain on large text corpora, fine-tune on specific tasks.
        \item In-context learning: LLMs adapt to new tasks with few examples.
    \end{itemize}

    \vspace{0.2cm}
    \textbf{Examples:}
    \begin{itemize}
        \item GPT models (autoregressive): Generate coherent long-form text.
        \item BERT (bidirectional): Better at understanding, classification, NER.
        \item Transformer architecture: Parallel processing, self-attention over context.
    \end{itemize}

    \vspace{0.2cm}
    \textit{Tradeoff: Larger models are more capable but require more compute and data.}
\end{frame}

\begin{frame}{Summary and Next Steps}
    \textbf{What we covered:}
    \begin{itemize}
        \item Preprocessing: tokenization, normalization, stemming.
        \item BoW representation: efficient, simple, limited.
        \item Sentiment analysis: lexicon-based and learning-based approaches.
        \item Applications: chatbots, translation, speech, and more.
    \end{itemize}

    \vspace{0.2cm}
    \textbf{Key takeaway:} NLP is a spectrum from simple pattern-matching to complex neural models.

    \vspace{0.2cm}
    \textbf{Up next (Lecture 6):} Introduction to Neural Networks --- the engine behind modern NLP.
    \begin{itemize}
        \item Biological inspiration and the perceptron.
        \item Multilayer networks and backpropagation.
    \end{itemize}
\end{frame}

\end{document}
